\documentclass[10pt,a4paper,sans]{moderncv}
\moderncvstyle{classic}
\moderncvcolor{blue}
\usepackage[utf8]{inputenc}
\usepackage[scale=0.85]{geometry}
\usepackage{helvet}

\name{Alexandro}{Disla}
\title{Économiste Appliqué | Statisticien | Spécialiste en Suivi et Évaluation}
\address{14, rue Saint-Louis Jeanty, Route de Frère}{Port-au-Prince, Haïti}{}
\phone[mobile]{(+509) 4148-3700}
\phone[fixed]{(+509) 4630-7530}
\email{alexandrodisla@hotmail.com}
\extrainfo{Né le 10 juillet 1991 à Port-au-Prince --- Nationalité haïtienne --- GitHub : \url{https://github.com/AD0791}}

\begin{document}
\makecvtitle

\section{Profil Professionnel}
\cvitem{}{Économiste appliqué formé au CTPEA sur un cycle de quatre ans en économie et statistique, avec plus de dix ans de pratique répartie entre l'\textbf{administration publique haïtienne} et les \textbf{programmes financés par des bailleurs internationaux}. Près de huit ans au Ministère de la Planification et de la Coopération Externe --- suivi de l'exécution du Plan Triennal d'Investissement, puis modélisation de prévision, de simulation et d'analyse d'impact des politiques publiques --- suivis de mandats de suivi-évaluation en santé, en eau et assainissement et en éducation. Je couvre la chaîne complète d'une évaluation : \textbf{conception} du dispositif de suivi et des indicateurs, \textbf{collecte} par questionnaires électroniques (ODK, KoboToolbox, CommCare), \textbf{traitement} et apurement des bases, \textbf{analyse} statistique et modélisation, puis \textbf{restitution} pour des lecteurs non techniciens.}

\section{Compétences Clés}
\cvitem{}{
\begin{itemize}
    \item \textbf{Conception d'indicateurs et de systèmes de suivi :} Cadres logiques, chaînes de résultats, tableaux de suivi des indicateurs (ITT) alignés sur les standards des bailleurs, matrices de résultats et indicateurs traceurs, valeurs de référence, cibles et désagrégations exigées.
    \item \textbf{Collecte et programmation d'enquêtes :} Questionnaires électroniques en XLSForm pour \textbf{ODK, KoboCollect et CommCare} --- logiques de saut, contraintes de validation, listes en cascade, formulaires multilingues, collecte hors ligne ; masques de saisie et contrôles quotidiens des enquêteurs.
    \item \textbf{Traitement, apurement et qualité des données :} Nettoyage, dédoublonnage, harmonisation et réconciliation de bases multi-sites ; protocoles de Data Quality Assessment établis et supervisés, avec remontée aux registres sources et suivi des actions correctives jusqu'à clôture.
    \item \textbf{Analyse statistique et modélisation :} Statistique descriptive et inférentielle, plans d'échantillonnage probabilistes et non probabilistes, pondération, désagrégation, analyse de tendances et d'écarts ; méthodes qualitatives et triangulation ; prévision, simulation et analyse d'impact de politiques publiques.
    \item \textbf{Coordination d'équipes, rédaction et restitution :} Supervision d'enquêteurs et d'agents de collecte multi-sites, encadrement technique de personnel de suivi-évaluation, modules de formation des équipes de collecte ; rapports de collecte et d'analyse, notes de politique publique, tableaux de bord et présentations PowerPoint de synthèse.
\end{itemize}}

\section{Suivi-Évaluation, Statistique et Politiques Publiques}

\cventry{Jan 2026 -- Mars 2026}{Data Analyst \& Consultant MEAL}{Anseye Pou Ayiti}{Télétravail}{}{
\begin{itemize}
    \item Définition de la stratégie de collecte et mise en œuvre d'indicateurs de performance pour des programmes de leadership multisectoriels.
    \item Codage des formulaires XLSForm pour ODK et Google Forms avec logiques de saut et contraintes de validation ; maintenance des scripts d'intégration Python sur AWS ; supervision de la qualité des données remontées et formation des agents de collecte.
\end{itemize}}

\cventry{Oct 2023 -- Jan 2026}{Consultant Suivi-Évaluation \& Analyste de Données}{HANWASH}{Télétravail}{}{
\begin{itemize}
    \item \textbf{Dispositif de suivi :} Raffinement des plans de suivi-évaluation et des tableaux de suivi des indicateurs (ITT) alignés sur les normes JMP pour des programmes d'eau et d'assainissement.
    \item \textbf{Enquêtes de terrain :} Déploiement d'enquêtes géospatiales complexes sur mWater, administration et apurement de la base, résolution des problèmes de collecte remontés du terrain.
    \item \textbf{Restitution et formation :} Tableaux de bord automatisés (mWater, Power BI) connectés aux API de collecte ; guides techniques et formation des équipes.
\end{itemize}}

\cventry{Oct 2021 -- Août 2024}{Officier de Suivi \& Évaluation}{Impact Youth Project, Caris Foundation International}{Haïti}{}{
\begin{itemize}
    \item \textbf{Systèmes de collecte :} Mise en place de systèmes de collecte numérique (CommCare, HIVHaiti) sur des programmes de santé multi-sites, avec standardisation des formulaires et des flux entre sites.
    \item \textbf{Assurance qualité (DQA) :} Établissement et supervision des protocoles de Data Quality Assessment sur des bases MySQL intégrées, avec remontée des chiffres rapportés aux registres sources.
    \item \textbf{Encadrement :} Supervision des agents de collecte sur l'ensemble des sites et encadrement technique du personnel de suivi-évaluation sur les procédures et les contrôles qualité.
    \item \textbf{Traitement et restitution :} Routines Python et SQL remplaçant des exports-imports manuels, réduisant le temps de traitement de 40 \% ; tableaux de bord (Power BI, Quarto) pour le suivi hebdomadaire des indicateurs.
\end{itemize}}

\cventry{Nov 2015 -- Août 2023}{Économiste-Planificateur --- Direction de Planification Économique et Sociale (DPES)}{Ministère de la Planification et de la Coopération Externe (MPCE)}{Port-au-Prince}{}{
\begin{itemize}
    \item \textbf{Suivi de programmes publics :} Contribution aux rapports de suivi de l'exécution du Plan Triennal d'Investissement 2014--2016 --- écarts entre programmation et réalisation, facteurs de retard.
    \item \textbf{Modélisation et analyse d'impact :} Affecté à l'unité statistique et modélisation, participation à la construction du << Modèle de Planification du Développement >>, outil de prévision, de simulation et d'analyse d'impact des politiques publiques, sous la supervision de Dr. Edouard Nsimba, expert international en modélisation.
    \item \textbf{Cadrage et rédaction :} Réunions de cadrage macroéconomique avec le Ministère de l'Économie et des Finances (MEF), suivi des indicateurs macroéconomiques et rédaction pour des rapports de politiques publiques.
\end{itemize}}

\section{Ingénierie de Données et Systèmes}

\cventry{Févr 2026 -- Présent}{Ingénieur Logiciel (Fullstack)}{Tekkod LLC}{Télétravail}{}{Responsable de l'analyse de données, de la visualisation, des tableaux de bord et de l'ingénierie de données au sein des projets.}

\cventry{Mars 2021 -- Mai 2024}{Développeur Backend \& Ingénieur de Données}{Tekkod LLC}{Télétravail}{}{Pipelines de traitement de données (ETL, Apache Beam) migrant des données MongoDB vers BigQuery ; administration de bases relationnelles en accès concurrent et optimisation des requêtes.}

\cventry{Nov 2015 -- Mars 2021}{Consultant Logiciel \& Spécialiste Intégration de Données}{CassionSoft (Projet UNOPS)}{Haïti}{}{Intégration de données sur un projet des Nations Unies : conception de schémas relationnels, scripts de migration et validation de la logique d'intégration par tests unitaires.}

\section{Formation}
\cvitem{2010--2014}{\textbf{Diplôme d'Études Supérieures en Économie Appliquée} (niveau licence, cycle de quatre ans en économie et statistique) --- C.T.P.E.A, Centre de Techniques de Planification et d'Économie Appliquée, établissement sous tutelle du MPCE, Port-au-Prince. \textit{Lauréat}.}
\cvitem{2009--2010}{Baccalauréat, Philo section C, mention Bien --- Institution Saint-Louis de Gonzague, Port-au-Prince}

\section{Compétences Techniques}
\cvitem{Statistique}{\textbf{Stata}, \textbf{R} (RMarkdown), \textbf{Python} (Pandas, NumPy), \textbf{SQL}, \textbf{Excel avancé}, Microsoft Access}
\cvitem{Collecte}{\textbf{ODK}, \textbf{KoboToolbox}, \textbf{CommCare}, XLSForm, mWater, HIVHaiti}
\cvitem{Restitution}{\textbf{Power BI} (expert), Tableau, Looker Studio, Quarto, PowerPoint ; bases MySQL, PostgreSQL, SQLite, MongoDB, BigQuery}
\cvitem{Langues}{Créole (maternel), Français (maternel), Anglais (lu et écrit : excellent ; parlé : moyen)}

\end{document}
